\section{Concluding remarks}
\label{sec:discussion}

The finite-window model singles out a concrete obstruction to restricting Fibonacci expansions before normalization. The relevant recovery thresholds separate cleanly. For \(m\ge3\), \(m\) consecutive fold labels recover the underlying block of length \(2m-1\), and this is sharp because collisions persist at length \(2m-2\). For the local inverse, however, three labels determine the current digit and two do not. Thus the exact zero-anticipation memory is \(2\), independently of \(m\).

The graph and decoding statements arise from distinct mechanisms. A zero-synchronized span-\(r\) normalizer satisfying the bounded-cascade hypothesis produces a synchronized pair shift with memory \(r-1\), but individual rewrite support does not imply that hypothesis. Fibonacci normalization has unbounded rewrite cascades; its four-state graph instead follows from the explicit carry imbalance. The boundary identity \eqref{eq:window-boundary-identity} separately controls both whole-block separation and the shorter causal decoder. Passing from the raw/normalized pair to its discrepancy bit produces a strictly sofic factor with entropy \(\log\varphi\), so the anomaly retains a complete golden-mean exponential scale while losing the finite-type property. The joint-density pentagon adds a rigidity statement not visible in either marginal: maximal discrepancy at fixed output density is supported on one of two mixing frontier systems, both governed by the plastic constant. The discrepancy laws, decoder theorem, and weighted-fiber asymptotics are therefore linked to explicit structural arguments rather than independent finite-state computations.
